\documentclass[12pt, a4paper]{article}

\usepackage{amsmath, amssymb, amsthm}
\usepackage{geometry}
\usepackage{hyperref}
\usepackage{booktabs}
\usepackage{array}
\usepackage{xcolor}
\usepackage{titlesec}
\usepackage{enumitem}
\usepackage{microtype}

\geometry{margin=2.5cm}

\hypersetup{
    colorlinks=true,
    linkcolor=blue!60!black,
    urlcolor=blue!60!black,
}

% Theorem environments
\theoremstyle{definition}
\newtheorem{definition}{Definition}[section]

\theoremstyle{plain}
\newtheorem{proposition}{Proposition}[section]
\newtheorem{corollary}{Corollary}[section]

\theoremstyle{remark}
\newtheorem{remark}{Remark}[section]

% Operators
\DeclareMathOperator*{\argmax}{arg\,max}
\newcommand{\R}{\mathbb{R}}
\newcommand{\Stwo}{S^2}
\newcommand{\norm}[1]{\left\lVert#1\right\rVert}
\newcommand{\dd}{\,\mathrm{d}}

\title{%
  \textbf{Mathematical Proof: Level-0 Spherical Harmonics}\\[4pt]
  \large as the Exact and Complete Representation\\
  for Isotropic Fluorescence Emission in our model
}
\author{Armin Sheibanifard\\
\small National Centre for Computer Animation, Bournemouth University}
\date{\today}

% ===================================================================
\begin{document}
\maketitle
\tableofcontents
\bigskip

% -------------------------------------------------------------------
\section{Introduction}

This document provides a rigorous mathematical proof of the following
central hypothesis:

\begin{quote}
\itshape
For isotropic fluorescence emission as modelled in NeuroSGM, the
spherical harmonic expansion truncated at degree $L=0$ is not merely
a simplification but the \emph{exact} and \emph{physically necessary}
representation.  Consequently, (i)~the $L=0$ basis function
$Y_0^0$ is a positive constant with zero oscillatory content, and
(ii)~no negative oscillation exists at this level, rendering concerns
about the negative portions of higher-degree sinusoidal harmonics
entirely vacuous.
\end{quote}

% -------------------------------------------------------------------
\section{Definitions}

\begin{definition}[Unit sphere]
\label{def:sphere}
The unit sphere in $\R^3$ is
\[
  \Stwo \;=\; \bigl\{\,\mathbf{d}\in\R^3 : \norm{\mathbf{d}}=1\,\bigr\}.
\]
Every direction $\mathbf{d}\in\Stwo$ is parameterised by polar angle
$\theta\in[0,\pi]$ and azimuthal angle $\phi\in[0,2\pi)$ via
\[
  \mathbf{d} = (\sin\theta\cos\phi,\;\sin\theta\sin\phi,\;\cos\theta)^\top.
\]
\end{definition}

\begin{definition}[Real spherical harmonics]
\label{def:sh}
For degree $l\in\mathbb{N}_0$ and order $m\in\{-l,\ldots,l\}$, the
\emph{real spherical harmonic} $Y_l^m : \Stwo\to\R$ is
\[
  Y_l^m(\theta,\phi)
  \;=\;
  N_l^m\;\cdot\;P_l^{|m|}(\cos\theta)\;\cdot\;
  \begin{cases}
    \sqrt{2}\,\cos(m\phi)   & m > 0,\\[2pt]
    1                        & m = 0,\\[2pt]
    \sqrt{2}\,\sin(|m|\phi) & m < 0,
  \end{cases}
\]
where
\[
  N_l^m \;=\; \sqrt{\frac{2l+1}{4\pi}\,\frac{(l-|m|)!}{(l+|m|)!}}
\]
is the normalisation constant and $P_l^{|m|}$ denotes the associated
Legendre polynomial of degree $l$ and order $|m|$.
\end{definition}

\begin{definition}[Orthonormality of SH]
\label{def:ortho}
The real spherical harmonics form an orthonormal basis for
$L^2(\Stwo)$:
\[
  \int_{\Stwo} Y_l^m(\mathbf{d})\;Y_{l'}^{m'}(\mathbf{d})\dd\mathbf{d}
  \;=\;
  \delta_{ll'}\,\delta_{mm'}.
\]
\end{definition}

\begin{definition}[SH series expansion]
\label{def:shexp}
Any square-integrable function $f:\Stwo\to\R$ admits the unique
decomposition
\[
  f(\mathbf{d})
  \;=\;
  \sum_{l=0}^{\infty}\sum_{m=-l}^{l}
    c_l^m\;Y_l^m(\mathbf{d}),
  \qquad
  c_l^m \;=\; \int_{\Stwo} f(\mathbf{d})\;Y_l^m(\mathbf{d})\dd\mathbf{d},
\]
where $c_l^m\in\R$ are the \emph{spectral coefficients}.
\end{definition}

\begin{definition}[Fluorescence intensity model]
\label{def:intensity}
Each 3D Gaussian primitive $i$ is assigned an \emph{emission intensity}
$a_i : \Stwo \to \R_{>0}$, mapping the unit viewing direction
$\mathbf{d}_i\in\Stwo$ to a non-negative scalar signal amplitude.
\end{definition}

\begin{definition}[MIP rendering operator]
\label{def:mip}
Let $K$ be the number of Gaussian primitives.  The rendered intensity
at image pixel $(u,v)$ under the \emph{Maximum Intensity Projection}
(MIP) operator is
\[
  I(u,v)
  \;=\;
  \max_{i=1}^{K}
    \Bigl[\,a_i(\mathbf{d}_i)\;\cdot\;G_{2D,i}(u,v)\,\Bigr],
\]
where
\begin{align*}
  G_{2D,i}(u,v)
  &\;=\;
  \exp\!\Bigl(
    -\tfrac{1}{2}\,
    (\mathbf{p}-\boldsymbol{\mu}_{2D,i})^\top
    \Sigma_{2D,i}^{-1}
    (\mathbf{p}-\boldsymbol{\mu}_{2D,i})
  \Bigr)
  \;\in\;(0,1],\\[4pt]
  \mathbf{d}_i
  &\;=\;
  \frac{\boldsymbol{\mu}_i^{\mathrm{cam}}}
       {\norm{\boldsymbol{\mu}_i^{\mathrm{cam}}}},
\end{align*}
$\boldsymbol{\mu}_{2D,i}$ is the projected 2D centre of Gaussian $i$,
$\Sigma_{2D,i}$ is its projected 2D covariance, and
$\boldsymbol{\mu}_i^{\mathrm{cam}}$ is the Gaussian centre in camera
coordinates.
\end{definition}

% -------------------------------------------------------------------
\section{The \texorpdfstring{$L=0$}{L=0} Spherical Harmonic}

\begin{proposition}
\label{prop:Y00constant}
The $L=0$ spherical harmonic $Y_0^0$ is a strictly positive constant
on $\Stwo$, independent of direction.
\end{proposition}

\begin{proof}
Setting $l=0$ and $m=0$ in Definition~\ref{def:sh}:
\[
  N_0^0
  \;=\;
  \sqrt{\frac{2\cdot 0+1}{4\pi}\,\frac{(0-0)!}{(0+0)!}}
  \;=\;
  \sqrt{\frac{1}{4\pi}}.
\]
The associated Legendre polynomial at $l=0$, $m=0$ evaluates to
\[
  P_0^0(\cos\theta) \;=\; 1 \qquad \forall\,\theta\in[0,\pi],
\]
and the azimuthal factor for $m=0$ is $1$ by Definition~\ref{def:sh}.
Therefore
\[
  \boxed{Y_0^0(\theta,\phi)
  \;=\;
  \sqrt{\frac{1}{4\pi}}\;\cdot\;1\;\cdot\;1
  \;=\;
  \frac{1}{2\sqrt{\pi}}
  \;\approx\;
  0.2821}
\]
for all $(\theta,\phi)$, i.e.\ for all $\mathbf{d}\in\Stwo$.  The
value is strictly positive and depends on neither $\theta$ nor $\phi$.
\end{proof}

% -------------------------------------------------------------------
\section{Direction-Independence at \texorpdfstring{$L=0$}{L=0}}

\begin{proposition}
\label{prop:dirindep}
Truncating the SH expansion at degree $L=0$ yields a
direction-independent function.
\end{proposition}

\begin{proof}
From Definition~\ref{def:shexp}, the $L=0$ truncation retains only
$l=0$, $m=0$:
\[
  f(\mathbf{d})
  \;\Big|_{L=0}
  \;=\;
  \sum_{m=-0}^{0} c_0^m\,Y_0^m(\mathbf{d})
  \;=\;
  c_0^0\,Y_0^0.
\]
By Proposition~\ref{prop:Y00constant}, $Y_0^0 = \frac{1}{2\sqrt{\pi}}$
is constant.  Hence
\[
  f(\mathbf{d})\Big|_{L=0}
  \;=\;
  \frac{c_0^0}{2\sqrt{\pi}}
  \;=\;
  \mathrm{const}
  \qquad
  \forall\,\mathbf{d}\in\Stwo.
\]
The function value is identical for every direction on the sphere.
\end{proof}

% -------------------------------------------------------------------
\section{Non-Negativity of the Emission Amplitude}

\begin{proposition}
\label{prop:nonneg}
Let $c_i\in\R$ be a learned scalar parameter for Gaussian $i$.  The
sigmoid-activated $L=0$ intensity
\[
  a_i(\mathbf{d})
  \;=\;
  \sigma\!\left(\frac{c_i}{2\sqrt{\pi}}\right)
  \;=\;
  \sigma\bigl(c_i\,Y_0^0\bigr)
\]
satisfies $a_i(\mathbf{d})\in(0,1)$ for all $\mathbf{d}\in\Stwo$ and
all $c_i\in\R$, where $\sigma(x)=\bigl(1+e^{-x}\bigr)^{-1}$.
\end{proposition}

\begin{proof}
Since $Y_0^0$ is constant (Proposition~\ref{prop:Y00constant}), the
argument of $\sigma$ is the fixed real scalar $c_i/(2\sqrt{\pi})$,
independent of $\mathbf{d}$.  The sigmoid satisfies
\[
  \sigma : \R \;\longrightarrow\; (0,1) \quad \text{strictly},
\]
because $e^{-x}>0$ for all $x\in\R$, giving $1+e^{-x}>1$ and thus
$\sigma(x)<1$, while $e^{-x}<\infty$ gives $\sigma(x)>0$.  Therefore
\[
  a_i(\mathbf{d})
  \;=\;
  \sigma\!\left(\frac{c_i}{2\sqrt{\pi}}\right)
  \;\in\;(0,1)
  \qquad
  \forall\,\mathbf{d}\in\Stwo,\;\forall\,c_i\in\R.
  \qedhere
\]
\end{proof}

% -------------------------------------------------------------------
\section{Absence of Oscillation at \texorpdfstring{$L=0$}{L=0}}

\begin{proposition}
\label{prop:nooscill}
$Y_0^0$ contains no oscillatory component.  Specifically, it has zero
spatial frequency in both $\theta$ and $\phi$.
\end{proposition}

\begin{proof}
From Proposition~\ref{prop:Y00constant}, $Y_0^0 = \tfrac{1}{2\sqrt{\pi}}$.
Taking partial derivatives:
\[
  \frac{\partial\,Y_0^0}{\partial\theta} \;=\; 0,
  \qquad
  \frac{\partial\,Y_0^0}{\partial\phi}   \;=\; 0.
\]
The function is flat on $\Stwo$.  In particular:
\begin{itemize}[leftmargin=2em]
  \item There is \emph{no} $\phi$-dependent factor $\cos(m\phi)$ or
        $\sin(|m|\phi)$, because the $m=0$ case in
        Definition~\ref{def:sh} contributes a factor of $1$.
  \item There is \emph{no} Legendre polynomial variation in $\theta$,
        because $P_0^0(\cos\theta)=1$ for all $\theta$.
\end{itemize}
Hence $Y_0^0$ oscillates with zero frequency in both angular
coordinates.  It attains the single value $\tfrac{1}{2\sqrt{\pi}}>0$
everywhere, so it has \emph{no negative values}.
\end{proof}

\begin{corollary}
\label{cor:noneg}
Concerns about the negative portions of higher-degree sinusoidal
harmonics are vacuous at $L=0$.  There exist no negative values to
suppress or discard.
\end{corollary}

\begin{proof}
Immediate from Proposition~\ref{prop:nooscill}: $Y_0^0>0$ everywhere.
\end{proof}

% -------------------------------------------------------------------
\section{Physical Justification: Isotropic Emission}

\begin{proposition}
\label{prop:physical}
For isotropic fluorescence emission, $L=0$ SH is the uniquely correct
and complete truncation level.
\end{proposition}

\begin{proof}
A fluorophore emits photons \emph{isotropically}: its emission
intensity is identical in every direction,
\[
  a_i(\mathbf{d}) \;=\; a_i \;=\; \mathrm{const}
  \qquad
  \forall\,\mathbf{d}\in\Stwo.
\]
By Definition~\ref{def:shexp}, the spectral coefficients of $a_i$ are
\[
  c_l^m
  \;=\;
  \int_{\Stwo} a_i\;Y_l^m(\mathbf{d})\dd\mathbf{d}
  \;=\;
  a_i \int_{\Stwo} Y_l^m(\mathbf{d})\dd\mathbf{d}.
\]
By orthonormality (Definition~\ref{def:ortho}), for $l\geq 1$:
\[
  \int_{\Stwo} Y_l^m(\mathbf{d})\dd\mathbf{d}
  \;=\;
  \int_{\Stwo} Y_l^m(\mathbf{d})\;Y_0^0\cdot\frac{1}{Y_0^0}\dd\mathbf{d}
  \;=\;
  \frac{1}{Y_0^0}
  \underbrace{\int_{\Stwo} Y_l^m(\mathbf{d})\;Y_0^0\dd\mathbf{d}}_{=\,\delta_{l0}\delta_{m0}\,=\,0}
  \;=\; 0.
\]
Therefore
\[
  c_l^m \;=\; 0
  \qquad
  \forall\, l\geq 1,\; m\in\{-l,\ldots,l\}.
\]
The \emph{exact} SH representation of $a_i$ is
\[
  a_i(\mathbf{d})
  \;=\;
  c_0^0\,Y_0^0
  \;=\;
  \frac{c_0^0}{2\sqrt{\pi}},
\]
containing \emph{only} the $l=0$ term.  Including any $l\geq 1$ terms
would model directional variation that is physically absent, introducing
unnecessary parameters without improving expressive power for this
signal class.
\end{proof}

% -------------------------------------------------------------------
\section{Equivalence to the Current Scalar Model}

\begin{proposition}
\label{prop:equiv}
The NeuroSGM scalar intensity model
$a_i = \sigma(\mathrm{log\_intensity}_i)$
is algebraically equivalent to the $L=0$ SH model.
\end{proposition}

\begin{proof}
In the $L=0$ SH model (Proposition~\ref{prop:dirindep}):
\[
  a_i
  \;=\;
  \sigma\!\left(c_0^0\cdot Y_0^0\right)
  \;=\;
  \sigma\!\left(\frac{c_0^0}{2\sqrt{\pi}}\right).
\]
Define $\tilde{c}_i \;=\; c_0^0/(2\sqrt{\pi}) \;\in\;\R$.  Then
\[
  a_i \;=\; \sigma(\tilde{c}_i).
\]
Setting $\tilde{c}_i = \mathrm{log\_intensity}_i$ recovers exactly
the current model.  The two parameterisations differ only by the
absorbing constant $1/(2\sqrt{\pi})$ into the learned scalar, which
does not affect expressiveness.
\end{proof}

% -------------------------------------------------------------------
\section{Summary}

The following chain of implications constitutes the complete proof:

\[
\underbrace{a_i(\mathbf{d})=\mathrm{const}}_{\substack{\text{isotropic}\\\text{emission}}}
\;\Longrightarrow\;
\underbrace{c_l^m=0\;\;\forall\,l\geq 1}_{\substack{\text{Proposition}\\\ref{prop:physical}}}
\;\Longrightarrow\;
\underbrace{Y_0^0=\tfrac{1}{2\sqrt{\pi}}}_{\substack{\text{Proposition}\\\ref{prop:Y00constant}}}
\;\Longrightarrow\;
\underbrace{\tfrac{\partial Y_0^0}{\partial\theta}=\tfrac{\partial Y_0^0}{\partial\phi}=0}_{\substack{\text{Proposition}\\\ref{prop:nooscill}}}
\;\Longrightarrow\;
\underbrace{a_i\in(0,1)}_{\substack{\text{Proposition}\\\ref{prop:nonneg}}}
\;\Longrightarrow\;
\underbrace{I(u,v)\in(0,1]}_{\substack{\text{Definition}\\\ref{def:mip}}}
\]

\bigskip

\begin{center}
\renewcommand{\arraystretch}{1.4}
\begin{tabular}{lcccc}
\toprule
\textbf{Property} & \textbf{Standard 3DGS} & \textbf{NeuroSGM ($L=0$)} \\
\midrule
Signal source           & External illumination   & Intrinsic emission \\
Direction dependence    & Yes (BRDF)              & \textbf{No (isotropic)} \\
SH degree required      & $L=1,2,3$               & \textbf{$L=0$ only} \\
$Y_l^m$ oscillates      & Yes ($l\geq 1$)         & \textbf{No ($l=0$)} \\
Negative SH values      & Present, managed        & \textbf{Non-existent} \\
Parameters per Gaussian & $(L+1)^2 \cdot 3$       & \textbf{1 scalar} \\
Physical correctness    & For reflectance         & \textbf{Exact for fluorescence} \\
\bottomrule
\end{tabular}
\end{center}

\bigskip
\noindent
\textbf{Conclusion.}\quad
The hypothesis is proved.  $L=0$ SH is not a simplification of a
richer directional model; it is the \emph{mathematically exact and
physically necessary} representation for isotropic fluorescence
emission.  The concern about negative oscillations is vacuous at
$L=0$ because $Y_0^0$ is a strictly positive constant with identically
zero spatial frequency in all directions.  The current NeuroSGM scalar
intensity model $a_i=\sigma(\mathrm{log\_intensity}_i)$ is the
complete and correct SH representation for this physical system.
\hfill$\blacksquare$

\end{document}
